\documentclass{report}
\usepackage{graphicx} % Required for inserting images

\usepackage{fontspec}
\setmainfont{FreeSerif}

\usepackage{polyglossia}
\setdefaultlanguage{greek}
\setotherlanguages{english}
\usepackage[normalem]{ulem}

\usepackage{caption}
\usepackage[useregional]{datetime2}
\usepackage{amsmath}

\usepackage{url}
\usepackage{hyperref}
\usepackage{biblatex}
\bibliography{citation}
%\addbibresource{citation}
\usepackage[dvipsnames]{xcolor}

\hypersetup{
	colorlinks=true,
	linkcolor=black,
	filecolor=magenta,      
	urlcolor=cyan,
	citecolor = violet
}

\urlstyle{same}

%\DTMlangsetup{showdayofmonth=false}

\title{Προκαταρκτική Παραγωγή Καμπυλών Επικινδυνότητας στο Μετρό Θεσσαλονίκης}
\author{Διαμαντής Χατζηαργυρίου}
\date{\Today}

\begin{document}

\maketitle

\section{Εισαγωγή}

Πριν λίγες εβδομάδες, ξεκίνησα να χρησιμοποιώ το Open Quake Engine για την ποσοτικοποίηση της επικινδυνότητας στο μετρό Θεσσαλονίκης.

Μετά από μελέτη των σχετικών οδηγιών του προγράμματος και εκτέλεση μερικών αναλύσεων, συντάσσεται η συγκεκριμένη αναφορά για οργάνωση προκαταρκτικών συμπερασμάτων όπως συμπυκνώνονται στο Σχήμα \ref{fig:haz}.

\section{Δεδομένα}
\label{data}
Για την ανάλυση, αρχικά, εξήχθη η χάραξη της πρώτης γραμμής του μετρό Θεσσαλονίκης από το OpenStreetMap.

Έπειτα, τα δεδομένα των συντεταγμένων τροχιών και σταθμών εμπλουτίστηκε μέσω του προγράμματος esrm20\_sitemodel με τιμές $V_s{30}$. Οι τιμές αυτές θα χρησιμεύσουν για την ακριβή εκτίμηση της επικινδυνότηταςσε κάθε σημείο. Έτσι, κατασκευάζονται τα site models που χρησιμοποιούνται για την ανάλυση της επικινδυνότητας.

Ακόμα, χρησιμοποιήθηκαν έτοιμα μοντέλα σεισμικών πηγών για την παραγωγή των κινήσεων.

Όλες οι αναλύσεις αναλύσης έτρεξαν με το 23ο seed, και 100 ή 1000 samples. Οι αναλύσεις έτρεξαν σε ικανοποιητικό χρόνο, οπότε μπορούν να αυξηθούν στο μέλλον.

\section{Ανάλυση}
Στην Ενότητα \ref{data}, έγινε αναφορά στα δεδομένα που συγκεντρώθηκαν για την ανάλυση.

Σε επόμενη φάση, τα συγκεκριμένα δεδομένα χρησιμοποιήθηκαν στο OQ ώστε να γίνει παραγωγή των αποτελεσμάτων.

Οι αναλύσεις περιλαμβάνουν αρχικά τον υπολογισμό της μέσης τιμής PGA κατά μήκος της δεύτερης τροχιάς, της πρώτης γραμμής, για έτη επαναφοράς 125, 600 και 2500. 

Τα έτη αυτά, επιλέχθησαν σύμφωνα με τις οριακές καταστάσεις DL, NC και SD αντίστοιχα, για κλάση επιπτώσεων 3, του αναθεωρημένου Ευρωκώδικα.

Έπειτα, πραγματοποιήθηκ ανάλυση στον σταθμό του πανεπιστημίου για την εξαγωγή καμπύλων επικινδυνότητας για μέγεθος σεισμικής έντασης PGA. Τρεις καμπύλες παρήχθησαν για την επιφάνεια και βάθη 10 και 40 μέτρων αντίστοιχα.

Ο σταθμός του πανεπιστημίου επιλέχθηκε καθώς είναι αντιπροσωπευτικό παράδειγμα των βαθών.

\section{Αποτελέσματα \& Συμπεράσματα}
Στο Σχήμα \ref{fig:haz} απεικονίζονται τα αποτελέσματα της μελέτης.

\subsection{Γραμμή 1 Τροχιά 2}
Πρώτον, αναπαριστώνται οι μέσες τιμές PGA σε σχέση με το μήκος της τροχιάς καθώς και της ταχύτητας μετάδοσης διατμητικών κυμάτων $V_s30$.

Από το σχήμα, φαίνεται ότι στα χαμηλές περιόδους επαναφοράς (125, 600), η διαμήκης διακύμανση της PGA είνα σχετικά χαμηλή, και εύκολα θα μπορούσε να παρεμβληθεί από μία ευθεία.

Σε περίοδο επαναφοράς 2500 ετών, υπάρχει μία σημαντικότερη διακύμανση, η οποία, όμως, κυμαίνεται περί των $0.75-0.9g$.

Συγκρίνοντας με τις τιμές $V_s30$, υποννοείται ελαφρώς μία ενίσχυση με την μείωση της τιμής, αλλά η συσχέτιση δεν είναι σαφής στην συγκεκριμένη ανάλυση και είναι προφανές ότι οι κύριες παράμετροι που την ορίζουν είναι η θέση μελέτης, η περίοδος επαναφοράς και οι πηγές διέγερσης.

\begin{figure}[h]
	\centering
	\includegraphics[width=0.85\linewidth]{../data/Plot1.png}
	\caption{Μέση PGA Γραμμής και Καμπύλες Επικινδυνότητας Μετρό Θεσσαλονίκης}
	\label{fig:haz}
\end{figure}

\subsection{Σταθμός Πανεπιστημίου}
Τέλος, παρήχθησαν τρεις καμπύλες επικινδυνότητας για διάφορα βάθη στην ίδια θέση.

Αυτές οι καμπύλες δίνουν μία τάξη μεγέθους της επικινδυνότητας στην ορισμένη θέση όπου επιτάχυνση $0.3g$, όπου τέτοια έργα -τυπικά- διαστασιολογούνται για να αντέχουν, ώστε να συμβαίνει μία ανά 100 έτη, και η μέγιστη επιτάχυνση φτάνει περί τα $0.9g$ στα 1000 έτη.

Η συσχέτιση με το βάθος φαίνεται χαμηλή, καθώς οι καμπύλες σχεδόν ταυτίζονται από $PoE = 10\%$ και έπειτα, όπου και συγκεντρώνεται η περισσότερη εκτόνωση σεισμικής ενέργειας.

\subsection{Συσχέτιση}
Προς το παρόν τα δεδομένα της $V_s30$ προέκυψαν αυτοματοποιημένα και όχι από μετρήσεις, και τα βάθη των σηράγγων ώστε να βελτιωθεί η ακρίβεια των αποτελεσμάτων.

Από τις αναλύσεις η συσχέτιση ίσως να μην είναι τόσο σημαντική, αλλά και πάλι θα χρειαστεί να συμπληρωθούν για πληρότητα.

\section{Επόμενα Βήματα}
Ολοκληρώνοντας τις αναλύσεις, διαφαίνονται κάποιες σαφείς κινήσεις που πρέπει να ακολουθήσουν, και αιωρούνται άλλες ενδεχόμενες.

Αρχικά, οι αναλύσεις πρέπει να επεκταθούν σε όλο το εύρος της γραμμής, ώστε να υπάρχει μία ολοκληρωμένη εικόνα.

Δηλαδή, θα πρέπει να συγκεντρωθούν τα δεδομένα της επέκτασης της Μίκρας, αλλά και των μελλοντικών επεκτάσεων. Αντίστοιχα χρειάζονται οι αυθεντικές χαράξεις, οι οποίες θα εμπεριέχουν τα βάθη των σηράγγων, αλλά και της ακριβείς θέσεις, καθώς το OSM δεν έχει ικανοποιητική ευκρίνεια για να βγάλει μία οιονεί συνεχή εικόνα της γραμμής. Αυτό το δεύτερο, διαφαίνεται και από το γεγονός ότι το συνολικό μήκος του συρμού υπολογίστηκε στα $8.4km$ σε αντίθεση με τα επίσημα $9.6km$.

Έπειτα, χρειάζεται να συγκεντρωθούν εδαφικά στοιχεία, για ακριβέστερο υπολογισμό της ταχύτητας $V_s30$, αλλά και πιθανής ανάλυσης επιρροής των τοπικών εδαφικών συνθηκών. Με ένα κατάλληλο προσομοίωμα θα μπορούσε ενδεχομένως να γίνει μία τέτοια ανάλυση στο Strata που θα ανήγαγε τις παραγώμενες καμπύλες σύμφωνα με κάποιες χαρακτηριστικές θέσεις.

Πρωτού αναλυθούν οι στρώσεις στο Strata, θα μπορούσε να εκτελεστεί μία ανάλυση με μία μέση $V_s30$ ώστε να αναλυθούν όλοι οι σταθμοί και να γίνει σύγκριση με τις αρχικές τιμές. Έτσι, θα διακρίνονταν αν όντως η επιρροή της $V_s30$ είναι χαμηλή, ή αν σφάλω.

Ακόμα, δεν γνωρίζω κατά πόσο συνηθίζεται, ή θα ήτα συνετό, αλλά σκοπεύω να παράξω μία τρισδιάστατη απεικόνιση μιας επιφάνειας επικινδυνότητας κατά μήκος του μετρό (ένωση των σημειακών καμπύλων επικινδυνότητας). Απλώς δεν ξέρω αν μπορεί να χρησιμοποιηθεί ως δεδομένο ακόμα, ή είναι φαμφάρα.

Ολοκληρώνοντας, αξίζει μία παρατήρηση -και ίσως η σημαντικότερη- ότι το Μέγεθος Έντασης που επιλέχθηκε ήταν η PGA. Αυτό δεν έγινε για κάποιον συγκεκριμένο σκοπό, απλώς για να παραχθούν κάποια πρωταρχικά αποτελέσματα, αλλά για σήραγγες μεσαίου προς μεγάλου βάθους, το κυρίως Μέγεθος Έντασης είναι η PGV. Επομένως, σε επερχόμενη φάση οι αναλύσεις θα πραγματοποιηθούν και με PGV.

\end{document}
